\scanpage{14}

b) $\forall\,F$ injectif sur $\unsure{Y}$, $g^{*}(F)$ est injectif et $F\unsure{\mathord{-}|\to}g_{*}g^{*}F$ est un épimorphisme (ex: $g$ inclusion d'un ouvert $U\hookrightarrow e_{Y}$).

On prend
\[
\mathbb R\Gamma_{Y}\bigl(\operatorname{Ker}\bigl(\underbrace{C(F}_{\text{rés. inj.}})\to g_{*}g^{*}\bigl(C(F)\bigr)\bigr)\bigr)
\]

Dans le cas général, la difficulté provient du fait que le cône d'un morphisme de complexes (tel que
\[
F\to g_{*}\deleted{\unsure{C|(}}\bigl(g^{*}(F)\deleted{\unsure{)|C}}\bigr)\deleted{\unsure{g_{*}|*}}\qquad\text{)}
\]
n'est pas fonctoriel (dans la catégorie dérivée) p. r. à la flèche dont on veut prendre le cône.

Et pourtant, dans le cas particulier actuel il devrait y avoir un choix fonctoriel. \unsure{ESt|Est}-ce \unsure{évident|evident} ?

\uline{Question pour Illusie} : Dans sa théorie des \unsure{déformations|déform.tions} de schémas en \unsure{groupes|groupe} plats, il tombe sur des $H^{3}(B_{G}/X,\quad)$ resp. des $\operatorname{Ext}^{2}(X;,\quad)$. Peut-on courtcircuiter sa théorie via la théorie (supposée écrite) des \unsure{Gr|gr}-champs -- resp. via la théorie des champs de Picard ? J'ai
% archive sheet number: 13
